% Part: sets-functions-relations
% Chapter: sets
% Section: proofs-about-sets

\documentclass[../../../include/open-logic-section]{subfiles}

\begin{document}

\olfileid[id]{sfr}{set}{prf}
\olsection{Bukti tentang Himpunan}

\begin{editorial}
Bagian ini digantikan oleh \verb|content/methods/proofs| dan secara bawaan
tidak disertakan dalam bab ini.
\end{editorial}

\begin{explain}
Himpunan dan notasi yang telah kita perkenalkan sejauh ini menyediakan
singkatan yang memudahkan kita menentukan himpunan dan menyatakan hubungan
di antaranya. Sering kali kita juga perlu membuktikan klaim tentang
hubungan tersebut. Jika Anda belum mengenal bukti matematika, hal ini
mungkin baru bagi Anda. Karena itu, kita akan membahas sebuah contoh
sederhana langkah demi langkah. Kita akan membuktikan bahwa untuk sebarang
himpunan $X$ dan $Y$, selalu berlaku $X \cap (X \cup Y) = X$. Bagaimana
cara membuktikan identitas antara himpunan seperti ini? Ingat bahwa
himpunan ditentukan semata-mata oleh !!{element}snya, yakni dua himpunan
identik jika dan hanya jika keduanya mempunyai !!{element}s yang sama.
Jadi, dalam hal ini kita harus membuktikan bahwa (a) setiap !!{element}
dari $X \cap (X \cup Y)$ juga merupakan !!a{element} dari $X$ dan,
sebaliknya, bahwa (b) setiap !!{element} dari $X$ juga merupakan
!!a{element} dari $X \cap (X \cup Y)$. Dengan kata lain, kita menunjukkan
baik (a) $X \cap (X \cup Y) \subseteq X$ maupun (b) $X \subseteq X \cap
(X \cup Y)$.

Klaim umum seperti ``setiap !!{element}~$z$ dari $X \cap (X \cup Y)$ juga
merupakan !!a{element} dari $X$'' dibuktikan dengan terlebih dahulu
mengasumsikan bahwa diberikan suatu $z \in X \cap (X \cup Y)$ yang
sebarang, lalu membuktikan dari asumsi ini bahwa $z \in X$. Anda mungkin
mengenal pola ini sebagai ``pembuktian bersyarat umum.'' Dalam bukti ini
kita juga harus menggunakan definisi-definisi yang muncul dalam asumsi dan
kesimpulan, misalnya, dalam hal ini, definisi ``$\cap$'' dan ``$\cup$.''
Jadi, kasus (a) dapat dikemukakan sebagai berikut:

\begin{quote}
(a) Pertama-tama kita ingin menunjukkan bahwa $X \cap (X \cup Y) \subseteq
X$, yakni, menurut definisi $\subseteq$, bahwa jika $z \in X \cap (X \cup
Y)$, maka $z \in X$, untuk sebarang~$z$. Jadi, asumsikan bahwa $z \in X
\cap (X \cup Y)$. Karena $z$ merupakan !!a{element} dari irisan dua
himpunan jika dan hanya jika ia merupakan !!a{element} dari kedua himpunan
tersebut, kita dapat menyimpulkan bahwa $z \in X$ dan juga $z \in X \cup
Y$. Khususnya, $z \in X$, dan itulah yang ingin kita tunjukkan.
\end{quote}

Ini menuntaskan paruh pertama bukti. Perhatikan bahwa pada langkah terakhir
kita menggunakan fakta bahwa jika suatu konjungsi ($z \in X$ dan $z \in X
\cup Y$) mengikuti dari sebuah asumsi, maka setiap konjungnya mengikuti
dari asumsi yang sama. Anda mungkin mengenal aturan ini sebagai
``eliminasi konjungsi,'' atau $\land$Elim. Sekarang mari kita buktikan (b):

\begin{quote}
(b) Sekarang kita membuktikan bahwa $X \subseteq X \cap (X \cup Y)$,
yakni, menurut definisi $\subseteq$, bahwa jika $z \in X$, maka $z \in X
\cap (X \cup Y)$ juga, untuk sebarang~$z$. Asumsikan $z \in X$. Untuk
menunjukkan bahwa $z \in X \cap (X \cup Y)$, kita harus menunjukkan
(menurut definisi ``$\cap$'') bahwa (i) $z \in X$ dan juga (ii) $z \in X
\cup Y$. Di sini, (i) hanyalah asumsi kita, sehingga tidak ada lagi yang
perlu dibuktikan. Untuk (ii), ingat bahwa $z$ merupakan !!a{element} dari
gabungan himpunan jika dan hanya jika ia merupakan !!a{element} dari
sekurang-kurangnya salah satu himpunan tersebut. Karena $z \in X$, dan $X
\cup Y$ adalah gabungan $X$ dan $Y$, hal ini berlaku di sini. Jadi, $z \in
X \cup Y$. Kita telah menunjukkan baik (i) $z \in X$ maupun (ii) $z \in X
\cup Y$; maka, menurut definisi ``$\cap$,'' $z \in X \cap (X \cup Y)$.
\end{quote}

Uraian ini agak panjang lebar, tetapi menggambarkan cara kita bernalar
tentang himpunan dan hubungan di antaranya. Biasanya kita tidak menjelaskan
semuanya secara sejelas ini; khususnya, kita mungkin tidak mengulang semua
definisi. Dalam teks yang lebih lanjut, bukti untuk hasil kita akan jauh
lebih padat. Bukti itu mungkin tampak seperti berikut ini.
\end{explain}


\begin{prop}[Absorpsi]
Untuk semua himpunan $X$, $Y$,
\[
X \cap (X \cup Y) = X
\]
\end{prop}

\begin{proof}
(a) Andaikan $z \in X \cap (X \cup Y)$. Maka $z \in X$, sehingga $X \cap
(X \cup Y) \subseteq X$.

(b) Sekarang andaikan $z \in X$. Maka $z \in X \cup Y$ juga, dan karena
itu $z \in X \cap (X \cup Y)$ juga. Dengan demikian, $X \subseteq X \cap
(X \cup Y)$.
\end{proof}

\begin{prob}
Buktikan secara terperinci bahwa $X \cup (X \cap Y) = X$. Kemudian berikan
versi bukti yang singkat dan padat. (Petunjuk: untuk arah $X \cup (X \cap Y)
\subseteq X$, Anda akan memerlukan pembuktian dengan kasus, yang juga
dikenal sebagai $\lor$Elim.)
\end{prob}

\end{document}
